\section{Números Reales}
\subsection{Representación digital del mundo real}
\begin{frame}{Números Reales}
 \begin{textblock*}{.5\textwidth}(0.4cm,1.2cm)
  \begin{itemize}[<+(1)->]\justifying  \parskip 0.2cm
   \item El rango de los números reales comprende desde $-\infty$ hasta $+\infty$.
   \item Los registros de un procesador tienen resolución finita.
   \item Por lo tanto un computador solo puede representar un sub conjunto de $\Re$.
   \item Además, no es solo un tema de magnitud sino de resolución.
  \end{itemize}
 \end{textblock*}
 \only<6->{
  \begin{textblock*}{.5\textwidth}((0.4cm+0.5\textwidth),1.2cm)
   \centering \includegraphics[width=\textwidth]{NumerosRealesenBinario.png}
  \end{textblock*}
 }
\end{frame}

% \subsection{Introducción}
\begin{frame}{Representación de Números Reales}
 \begin{textblock*}{\textwidth}(0.2cm,1.2cm)
  \begin{itemize} [<+(1)->]\justifying
   \item Un computador necesita soportar operaciones con números fraccionarios, denominados  en matemática como Números Reales.
   \item Ejemplos:
  \end{itemize}
 \end{textblock*}
 \begin{textblock*}{\textwidth}(0.9cm,3cm)
   \only<4->{$\pi=3.14159265358979323846..._{10}$\\}
   \only<5->{$e=2.7182818284590452353602874713527..._{10}$\\}
   \only<6->{$0.000000001_{10}$, ó $1.0_{10} x 10^{9}$ (nano metros por ejemplo)\\}
   \only<7->{$3,155,760,000$, ó $3.15576_{10} x 10^{9} $ segundos en un siglo\\}
 \end{textblock*}
 \begin{textblock*}{\textwidth}(0.2cm,4.9cm)
  \begin{itemize} [<+(5)->]\justifying
   \item En los tres primeros ejemplos se trata de números fraccionarios pequeños
   \item Pero el último es un número entero, pero tan grande que no puede ser representado mediante una variable entera de \SI{32}{\bit}
   \item En el tercer número se emplea \textbf{\textit{notación científica}}.
   \item En particular, cuando solo tiene un dígito no nulo como parte entera, y el resto se desplazan hacia la parte fraccionaria, ajustando el exponente, se habla de \textbf{\textit{Notación Científica Normalizada.}}
  \end{itemize}
 \end{textblock*}
\end{frame}

\begin{frame}{Notación Científica}
 \begin{textblock*}{\textwidth}(0.4cm,1.2cm)
  \begin{itemize}[<+->] \justifying
   \item Para el caso de los números reales se trabaja en \textbf{\textit{\color{blue}    notación científica}}.\\
   \begin{center}  \shadowbox{ $n= \pm f*10^e$ } \end{center}
   \item  $-725.832 = -7.25832*10^{2} = -725.832*10^{0}$
   \item  $3.14 = 0.314*10^{1} = 3.14*10^{0}$
   \item  $0.000001 = 0.1*10^{-5} = 1.0*10^{-6}$
   \item  $1941 = 0.1941*10^{4} = 1.941*10^{3}$
   \item Para unificar la representación  se recurre a la \textbf{\textit{\color{blue}notación científica normalizada}} \color{black}, en donde:\\
%   \vspace{0.3cm}
   \begin{center}
    $1 \leq$ \textbf{\textit{\color{blue}f}} \color{black} $<$ 9 ,y \textbf{\textit{\color{blue}e}} \color{black} es un entero con signo.
   \end{center}
  \end{itemize}
 \end{textblock*}
\end{frame}

\begin{frame}[t]{Notación Científica en el sistema Binario}
 \begin{textblock*}{\textwidth}(0.4cm,1.2cm)
  \begin{itemize} \justifying
   \item En el sistema binario la expresión de un número en notación científica normalizada es:\\
    \begin{center} \shadowbox{ $n= \pm f*2^e$  } \end{center}
   \vspace{0.3cm}
   \item En donde:\\
   \vspace{0.3cm}
   \begin{center}
 $0.5 \leq$ \textbf{\textit{\color{blue}f}} \color{black} $<$ 1 , 
y \textbf{\textit{\color{blue}e}} \color{black} es un entero con signo.
   \end{center}
  \end{itemize}
 \end{textblock*}
\end{frame}

\begin{frame}{Formatos de representación en un computador}
 \begin{textblock*}{\textwidth}(0.2cm,1.2cm)
  \begin{itemize} [<+(1)->]\justifying
   \item Al trabajar con números reales en un computador, enfrentamos un problema de primer magnitud: redondear cifras en una cantidad de bits acotada.
   \item Existen dos formatos de representación para número reales en un computador:
   \begin{itemize}
    \item Punto fijo
    \item Punto flotante
   \end{itemize}
   \item En general se impone la representación en punto flotante, dado que su rango de números representables con la misma cantidad de bits, es muchísimo mas amplio.
   \item La representación en Punto fijo permite utilizar los mismos registros de propósito general y la misma aritmética que se emplea con números enteros. La posición del punto decimal queda fija, por lo cual los resultados de las operaciones básicas son los mismos que para números enteros.
   \item Punto fijo es práctico en cálculos de rango numérico acotado. Ej: temperatura, presión, humedad, y demás magnitudes físicas cuyo rango es de algunos centenares de unidades a lo sumo, y requerimientos de precisión moderados.
  \end{itemize}
 \end{textblock*}
\end{frame}

\begin{frame}{Representación binaria en Punto Fijo con signo}
 \begin{textblock*}{\textwidth}(0.2cm,1.2cm)
  \begin{itemize}[<+(1)->]  \justifying \parskip 0.3cm
   \item Se representan mediante una expresión del tipo:\\
   $\color{blue}(a_{n}a_{n-1}\dots a_0.a_{-1}a_{-2}\dots a_{-m})_{2} = (-1)^{s}*(a_{n}*2^{n}+\dots + a_{0}*2^{0} + a_{-1}*2^{-1} + a_{-2}*2^{-2} +\dots+ a_{-m}*2^{-m})$
   \item Donde:
   \begin{itemize} \justifying
    \item \textbf{\textit{\color{blue}s}} \color{black}es el signo: \textbf{0} si el número es positivo y \textbf{1} si es negativo
    \item \textbf{\textit{\color{blue}$a_i$}} $\in  \mathbb{Z} $ y  $0 \leq$ \textbf{\textit{\color{blue}$a_i$}} $\leq 1$, $\forall$ \textbf{\textit{\color{blue}i}}\color{black} \textit{= -m, -1, 0, 1,\dots n}
   \end{itemize}
   \item Distancia entre dos números consecutivos es $2^{-m}$.
   \item \color{red}\textit{Deja de ser un rango continuo de números para transformarse en un rango discreto.}
  \end{itemize}
 \end{textblock*}
\end{frame}

\begin{frame}{Punto Flotante: Motivación}
 \begin{textblock*}{\textwidth}(0.2cm,1.2cm)
  \begin{itemize} [<+(1)->]\justifying
   \item Existe un rango amplio de aplicaciones que trabajan con números reales, y que además plantean requisitos muy estrictos de precisión.
   \item Para éstas resulta mucho mas adecuado el uso de punto flotante.
   \item Está basado en notación científica.
   \item En punto flotante un número se representa mediante un entero \textit{con signo} y de precisión fija llamado \textit{significando} o \textit{mantisa}, el que es escalado por otro entero también \textit{con signo} que se emplea como \textit{exponente} de una \textit{base} numérica fija.
   \item De este modo, cualquier número real puede ser representado en notación científica por medio de la siguiente expresión:
  \end{itemize}
  \only<6->{
   \begin{equation}\label{eq:scinot}
     n = m * b^e
   \end{equation}
   }
{\hangindent=0.8cm
\hangafter=1
\only<7->{\justifying donde \textit{m} representa la \textit{mantisa}, \textit{b} la \textit{base} numérica de representación, y \textit{e} el \textit{exponente} al que debe ser elevada la base.
}
\par}
 \end{textblock*}
\end{frame}

\begin{frame}{Punto Flotante: Precisión, Rango, Normalización}
 \begin{textblock*}{\textwidth}(0.2cm,1.2cm)
  \begin{itemize} [<+(1)->]\justifying
   \item La cantidad de bits de la \textit{mantisa} se relaciona con la precisión que puede obtenerse al representar un número irracional.
   \item Un número de precisión \textit{p} es aquel cuya representación emplea \textit{p} bits para la \textit{mantisa}, donde \textit{p} incluye el dígito correspondiente a su parte entera.
   \item El \textit{exponente}, por su parte, establece el rango de números que puede representarse.
   \item La expresión (1), sugiere que cada número puede tener tantas representaciones como combinaciones posibles existan entre los valores de \textit{exponente} y \textit{mantisa}.
   \item Para evitar tales ambig\"{u}edades, se recurre a una representación normalizada en que la \textit{mantisa} tiene un solo dígito entero. En \textit{base} 2, la simplificación evidente: hay solo dos valores posibles: $0$ o $1$.
   \item Si se asume que el dígito entero de la \textit{mantisa} siempre vale 1, se evita su inclusión.
   \item Se suma así a la \textit{mantisa} un bit adicional optimizando la precisión del número.
   \item Esta notación se denomina notación punto flotante normalizada.
  \end{itemize}
 \end{textblock*}
\end{frame}

\begin{frame}{Punto Flotante: Precisión, Rango, Normalización}
 \begin{textblock*}{\textwidth}(0.2cm,1.2cm)
  \begin{itemize} [<+(1)->]\justifying
   \item Un número real de precisión \textit{p} en base 2 puede expresarse de acuerdo con la siguiente expresión:\\
   \centering $n = \pm(1.m_{1}m_{2}m_{3}\dots m_{(p-2)}m_{(p-1)}) * 2^e$
   \item En su versión polinómica:\\
   \centering $n = \pm 1 + m_{1}\cdot 2^{-1} +  m_{2}\cdot 2^{-2} + \dot  m_{(p-1)}\cdot 2^{-(p-1)}$
   \item Esto sugiere que si aumentando la cantidad de bits del \textit{significando}, aumenta la precisión \textit{p}.
   \item Para una precisión \textit{p} dada, solo se puede representar un subconjunto de valores discretos pertenecientes a $\mathbb{R}$. Nunca al conjunto completo.
   \item En el siguiente slide representamos el conjunto discreto de Números Reales que pueden representarse para una dada Precisión \textit{p}.
   \item Puede verse que entre dos números consecutivos existe un intervalo de números reales que no pueden ser representados con esta Precisión ya que su \textit{mantisa} requiere más bits que la cantidad establecida por \textit{p}.
  \end{itemize}
 \end{textblock*}
\end{frame}

\begin{frame}{Punto Flotante: Precisión, Rango, Normalización}
 \begin{textblock*}{\textwidth}(0.2cm,1.2cm)
  \centering \includegraphics[width=.9\textwidth]{RealNumbers-1.pdf}
 \end{textblock*}
\end{frame}

\subsection{Normalización}
\begin{frame}{Situación inicial}
 \begin{textblock*}{\textwidth}(0.2cm,1.2cm)
  \begin{itemize} [<+->]\justifying \parskip .15cm
   \item Durante los años 60 y 70 cada fabricante de hardware que incluyó Unidades de Punto Flotante en la CPU de sus computadores optó por criterios de diseño particulares, en lo referente a la representación de los números.
   \item La cantidad de bits para \textit{mantisa}, \textit{exponente}, y la representación de cada uno de estos valores respecto del signo, así como los criterios para el redondeo en las operaciones, entre otros aspectos variaban de un procesador a otro.
   \item La precisión inherente de cada máquina era diferente, y por consecuencia, al ir acumulando resultados intermedios redondeados mediante criterios diferentes (truncado vs redondeo al mas próximo, por ejemplo), generaban resultados finales sensiblemente diferentes para un mismo algoritmo de cálculo (en rigor para algoritmos idénticos).
   \item Esta situación no resulta aceptable.
  \end{itemize}
 \end{textblock*}
\end{frame}

\begin{frame}{Situación inicial}
 \begin{textblock*}{\textwidth}(0.2cm,1.2cm)
  \begin{itemize} [<+->]\justifying \parskip .15cm
   \item Para solucionar esta divergencia se recurrió a una normalización del formato de representación de números en Punto Flotante.
   \item En 1985 se publica la primer especificación de una norma internacional a cargo del IEEE conocida como IEEE-754.
   \item Establece los formatos y criterios de cálculo para poner en común no solo el formato de representación, sino además otros aspectos relativos a las operaciones con estos números, tales como los criterios de redondeo, eventuales excepciones ante situaciones de borde como overflow, denormalización, etc.
   \item Por entonces, en 1985, la unidad de punto flotante que mas se había comercializado por lejos era el diseño de Intel de su Co Procesador de Punto Flotante 8087, presentado en 1978 junto al procesador 8086 y 8088 que IBM seleccionaría para su primer computador personal (IBM PC). Por tal motivo buena parte de la especificación IEEE-754 de 1985 ha tenido una marcada influencia de éste diseño.
  \end{itemize}
 \end{textblock*}
\end{frame}

\begin{frame}[t]{Formatos IEEE 754}
 \begin{textblock*}{\textwidth}(0.4cm,1.5cm)
  \only<1->{\includegraphics[width=\textwidth]{IEEE-754-layer1.pdf}}
 \end{textblock*} 
 \begin{textblock*}{\textwidth}(0.4cm,1.5cm)
  \only<2->{\includegraphics[width=\textwidth]{IEEE-754-layer2.pdf}}
 \end{textblock*} 
 \begin{textblock*}{\textwidth}(0.4cm,1.5cm)
  \only<3->{\includegraphics[width=\textwidth]{IEEE-754-layer3.pdf}}
 \end{textblock*} 
 \begin{textblock*}{\textwidth}(0.4cm,1.5cm)
  \only<4->{\includegraphics[width=\textwidth]{IEEE-754-layer4.pdf}}
 \end{textblock*} 
 \begin{textblock*}{\textwidth}(0.4cm,1.5cm)
  \only<5->{\includegraphics[width=\textwidth]{IEEE-754-layer5.pdf}}
 \end{textblock*} 
\end{frame}

\begin{frame}{IEEE-754}
 \begin{textblock*}{\textwidth}(0cm,1.2cm)
  \begin{itemize} [<+->]\justifying \parskip .12cm
   \item Cada formato consta finalmente de dos números de punto fijo: el \textit{significando}, y el \textit{exponente}. Ambos números signados.
   \item De acuerdo con la cantidad de bits de éstos, los números reales representables tienen diferentes precisiones, y rangos.
   \item Se han tomado algunas decisiones en la especificación del estándar, relativas a la representación de ambos números signados.
   \item El \textit{significando} se representa en signo y magnitud, y el \textit{exponente} en formato binario desplazado.
   \item La representación del \textit{exponente} tiene por objeto permitir la representación del número cero ($0_{10}$), ya que combinado con una mantisa nula con su parte entera también nula el resultado es cero, y todos los bits del número valen 0.
   \item Luego, cualquier otra combinación con \textit{exponente} $0$, pero con un \textit{significando} que tenga algún bit distinto de $0$ corresponderá a un número denormalizado no nulo.
  \end{itemize}
 \end{textblock*}
\end{frame}

\begin{frame}{IEEE-754}
 \begin{textblock*}{\textwidth}(0cm,1.2cm)
  \begin{itemize} [<+->]\justifying \parskip .15cm
   \item El cero tiene otras particularidades como la posibilidad de tener signo, es decir puede tenerse un valor $-0$ o $+0$. Desde el punto de vista operativo indicaría desde que extremo del conjunto de números se viene aproximando a cero un resultado, y permite además calcular la recíproca de ese numero para obtener $-\infty$.
   \item El bit mas significativo de cada representación se destina al signo de la \textit{mantisa}, o lo que es lo mismo, al signo del número representado.
   \item Continúa el \textit{exponente} en notación exceso (bias) o notación binario desplazado. y el resto de los bits representan el valor absoluto del \textit{significando}
   \begin{itemize}
    \item \textbf{\textit{Simgle Precision}}. Utiliza \SI{32}{\bit}, corresponde al tipo de dato \textit{float} del lenguaje C.
    \item \textbf{\textit{Double Precision}}. Utiliza \SI{64}{\bit}, correspondeal tipo de dato \textit{double} del lenguaje C.
    \item \textit{\textbf{Extended Double Precision}}. Utiliza \SI{80}{\bit}, derivado del diseño del Coprocesador Matemático que Intel había diseñado para su arquitectura x86 y que al momento del estándar era el mas difundido en el mercado.
   \end{itemize}
  \end{itemize}
 \end{textblock*}
\end{frame}


\begin{frame}{IEEE-754}
 \begin{textblock*}{\textwidth}(0cm,1.2cm)
  \begin{itemize} [<+->]\justifying \parskip .15cm
   \item Los formatos definidos en esta versión del estándar IEEE-754 junto con numerosas secciones, tales como redondeo, excepciones relacionadas con el resultado de acuerdo a condiciones particulares, entre otras reflejan una muy fuerte influencia de la arquitectura x86 de Intel en esta norma.
   \item Tal vez el formato Extended Double con el que en C se maneja el tipo de datos \textit{long double}, sea el signo mas notorio de esta influencia, ya que otros fabricantes de procesadores de la época como Motorola, manejaban el tipo \textit{long double} con valores de \SI{96}{\bit}.
   \item Por tal motivo el Extended Double, puede encontrarse normalmente como x86 Extended, o x86 \SI{80}{\bit}.
   \item Al corresponder Single Precision y Double Precision respectivamente a los tipos de datos \textit{float} y \textit{double} definidos en el Lenguaje de Programación C, se incluyó en IEEE-754 buena parte de las definiciones de ANSI respecto de estos formatos, y el ANSI participó muy activamente en el Comité que elaboró el estándar.
  \end{itemize}
 \end{textblock*}
\end{frame}

\begin{frame}{IEEE-754-2008}
 \begin{textblock*}{\textwidth}(0cm,1.2cm)
  \begin{itemize} [<+->]\justifying \parskip .15cm
   \item En 2008 se revisó la norma introduciéndose un formato corto de punto flotante (el denominado Half Precision cuyo tamaño es \SI{16}{\bit}), pero además se reemplaza el formato Extended Double Precision, por otro de \SI{128}{\bit} de ancho cuyo layout se muestra a continuación, y que para distinguirlo y ser breves denominamos Quad.
   \item Estos Formatos de Punto flotante son los establecidos en la revisión de la Norma IEEE754 en 2008.
   \item Single Precision y Double Precision  no cambian ya que son ampliamente utilizados. El formato Extended Double, por su parte plantea un datos de \SI{128}{\bit}, extendiendo básicamente la mantisa (precisión) en tanto el \textit{exponente} (rango) se mantiene. Y se agrega el formato Half precision.
   \item Puede inferirse que este cambio procura mayor precisión ya que la diferencia respecto del \textit{extended double} de 1985 está en la cantidad de bits del \textit{significando} que pasan de 64 a 112. El \textit{exponente} queda igual, de modo que el \textit{rango} de números reales representables no cambia.
  \end{itemize}
 \end{textblock*}
\end{frame}

\begin{frame}{IEEE-754-2008}
 \begin{textblock*}{\textwidth}(0cm,1.2cm)
  \centering \includegraphics[width=\textwidth]{IEEE754-2008.pdf}
 \end{textblock*}
\end{frame}

